% ===== PRÉAMBULE CANONIQUE — à coller VERBATIM en tête de CHAQUE .tex =====
\documentclass[11pt,a4paper]{article}
\usepackage[utf8]{inputenc}\usepackage[T1]{fontenc}\usepackage{lmodern}
\usepackage[french]{babel}
\usepackage{amsmath,amssymb,mathtools}\usepackage{icomma}
\usepackage[version=4]{mhchem}          % \ce{...} : équations & formules
\usepackage{siunitx}\sisetup{locale=FR} % \qty{}{}, \si{}, \ang{}
\usepackage{chemfig}                    % \chemfig{...} : structures développées
\usepackage{xcolor}\usepackage{colortbl}
\usepackage{tikz}\usepackage{pgfplots}\pgfplotsset{compat=1.18}
\usepackage[most]{tcolorbox}\usepackage{enumitem}\usepackage{array,booktabs}
\usepackage[a4paper,margin=2.2cm]{geometry}
\usetikzlibrary{arrows.meta,calc,angles,quotes,positioning,patterns,backgrounds,shapes.geometric}
\definecolor{primary}{RGB}{222,49,99}\definecolor{primarydark}{RGB}{150,22,55}
\definecolor{defcol}{RGB}{0,114,178}\definecolor{methcol}{RGB}{52,92,148}
\definecolor{excol}{RGB}{0,135,90}\definecolor{attncol}{RGB}{200,40,55}
\tcbset{boxbase/.style={enhanced,breakable,boxrule=0.9pt,arc=2.5pt,
  left=3mm,right=3mm,top=2.3mm,bottom=2.3mm,fonttitle=\bfseries,coltitle=white}}
% --- boîtes TUTORIEL (noms EXACTS, ne pas renommer) ---
\newtcolorbox{cle}[1][]{boxbase,colback=defcol!6,colframe=defcol,
  title={\textsc{À retenir}\ifx\relax#1\relax\relax\else\textmd{ : \itshape #1}\fi}}
\newtcolorbox{methode}[1][]{boxbase,colback=methcol!7,colframe=methcol,sharp corners,
  title={\textsc{Méthode}\ifx\relax#1\relax\relax\else\textmd{ --- \itshape #1}\fi}}
\newtcolorbox{exemplebox}[1][]{boxbase,colback=excol!5,colframe=excol,
  title={\textsc{Exemple}\ifx\relax#1\relax\relax\else\textmd{ : \itshape #1}\fi}}
\newtcolorbox{piege}[1][]{boxbase,colback=attncol!6,colframe=attncol,
  title={\textsc{Piège}\ifx\relax#1\relax\relax\else\textmd{ : \itshape #1}\fi}}
\newtcolorbox{remarque}[1][]{boxbase,colback=black!4,colframe=black!45,
  title={\textsc{Remarque}\ifx\relax#1\relax\relax\else\textmd{ : \itshape #1}\fi}}
% --- boîtes EXERCICE / CORRIGÉ (noms EXACTS, auto-comptées) ---
\newtcolorbox[auto counter]{exo}[1][]{boxbase,colback=primary!4,colframe=primary,
  title={\textsc{Exercice}~\thetcbcounter\ifx\relax#1\relax\relax\else\textmd{ --- \itshape #1}\fi}}
\newtcolorbox[auto counter]{corrige}[1][]{boxbase,colback=excol!4,colframe=excol,
  title={\textsc{Corrigé}~\thetcbcounter\ifx\relax#1\relax\relax\else\textmd{ --- \itshape #1}\fi}}
\newcommand{\R}{\mathbb{R}}\newcommand{\N}{\mathbb{N}}\newcommand{\Z}{\mathbb{Z}}\newcommand{\ds}{\displaystyle}
% (un \usepackage{fancyhdr}+en-tête est possible ; il est ignoré côté web)
% =========================================================================
\begin{document}

\begin{exo}[deux ions de valence 1]
Quand cation et anion sont tous deux de valence 1, la formule est en rapport $1:1$. Écris la formule brute
neutre.
\begin{enumerate}[label=\alph*)]
  \item \ce{Na+} et \ce{Cl-}
  \item \ce{K+} et \ce{NO3-}
  \item \ce{NH4+} et \ce{Cl-}
  \item \ce{Ag+} et \ce{Br-}
\end{enumerate}
\end{exo}

\begin{exo}[croiser une valence 1 et une valence 2]
Croise les valences pour équilibrer les charges. Écris la formule brute.
\begin{enumerate}[label=\alph*)]
  \item \ce{Ca^2+} et \ce{Cl-}
  \item \ce{Na+} et \ce{O^2-}
  \item \ce{K+} et \ce{S^2-}
  \item \ce{Al^3+} et \ce{Cl-}
\end{enumerate}
\end{exo}

\begin{exo}[penser à simplifier]
Ici le croisement donne des indices simplifiables : réduis-les au plus petit rapport entier.
\begin{enumerate}[label=\alph*)]
  \item \ce{Mg^2+} et \ce{O^2-}
  \item \ce{Ca^2+} et \ce{S^2-}
  \item \ce{Al^3+} et \ce{PO4^3-}
\end{enumerate}
\end{exo}

\begin{exo}[mettre l'ion polyatomique entre parenthèses]
L'anion est polyatomique : dès qu'il est présent au moins deux fois, on l'entoure de parenthèses.
\begin{enumerate}[label=\alph*)]
  \item \ce{Mg^2+} et \ce{OH-}
  \item \ce{Ca^2+} et \ce{NO3-}
  \item \ce{NH4+} et \ce{SO4^2-}
\end{enumerate}
\end{exo}

\begin{exo}[quand une valence 3 entre en jeu]
Croise en tenant compte de la valence 3. Écris la formule brute neutre.
\begin{enumerate}[label=\alph*)]
  \item \ce{Al^3+} et \ce{O^2-}
  \item \ce{Fe^3+} et \ce{Cl-}
  \item \ce{Ca^2+} et \ce{PO4^3-}
\end{enumerate}
\end{exo}

\end{document}
